\documentclass[UTF8,fontset=none,11pt,a4paper]{ctexart}
\usepackage{ipara-handbook}

\handbooksetup
  {局部到整体：中值定理与函数形状}
  {数学一 · 高等数学 Topic 04}
  {中值定理 · 零点 · 函数形状}

\begin{document}
\handbooktitle
  {局部到整体：中值定理与函数形状}
  {从“存在一个中值点”到“恢复整个区间的行为”}
  {}

\begin{quote}
本册只追问一个根本问题：\textbf{已知函数在端点、零点或局部导数上的信息，怎样把它运输成区间内部的见证点和整体形状？} 中值定理负责产生内部证据，辅助函数负责把目标表达式翻译成可用的导数结构，导数符号再把点态信息恢复成单调、极值、凹凸和零点分布。
\end{quote}

\begin{center}\rule{0.5\linewidth}{0.5pt}\end{center}

% ============================================================
\setcounter{section}{-1}
\section{本册定位}
% ============================================================

本册是高等数学 Subject Atlas 下的 Topic04。Topic02 提供连续与极限资格，Topic03 提供导数和有限局部模型；本册研究这些局部/边界信息怎样借助区间定理产生整体结论。

\subsection{Own}

本册独立负责：
\begin{itemize}
  \item 介值定理、零点定理、Rolle、Lagrange 与 Cauchy 中值定理的资格和逻辑链；
  \item 证明题中由目标表达式反推辅助函数的构造机制；
  \item 零点、重零点与反复 Rolle 的区间见证；
  \item 导数符号到单调性、极值、凹凸、拐点和函数形状的恢复；
  \item 驻点、临界点、极值点和拐点的候选/确认边界；
  \item 用零点重数、导数根的符号与 Rolle 链分析根的数量；
  \item 中值定理证明与函数形状题的统一检查协议。
\end{itemize}

\subsection{Uses / Stop Boundary}

本册调用 Topic01 的函数对象/定义域，Topic02 的连续性，Topic03 的局部导数和有限 Taylor。以下内容只保留接口：
\begin{itemize}
  \item 多中值点插点法、中值点参数渐近位置、插值余项与导数估计归 H-B02；
  \item Taylor 多项式本体归 Topic03；Lagrange 型余项的区间误差接口由 H-B02 接收；
  \item 变上限积分、原函数和积分因子所需的正则性归 Topic05/H-B03；
  \item 一阶微分方程的整体解族归 Topic12，本册只把积分因子作为辅助函数构造的局部工具；
  \item 单一问题族的题目信号、首步动作与退出条件进入同目录训练 Markdown；跨多个训练专题的稳定控制才进入《高等数学做题规则》。
\end{itemize}

% ============================================================
\section{根本问题：端点信息怎样变成区间内部证据？}
% ============================================================

只知道 $f(a)$、$f(b)$ 或 $f'(a)$，通常不能直接写出某个内部点 $\xi$。朴素做法是“找一个好看的点”或直接把平均斜率当作某点斜率，但这缺少存在性证明，也没有解释为什么目标表达式会在该点出现。

中值定理提供一条受资格条件保护的运输通道：
\begin{itemize}
  \item 连续性把端点值之间的中间输出全部填满；
  \item 端点等值把函数压成一条必须在内部水平的轨迹；
  \item 端点不等时，减去割线后制造端点等值，得到平均变化率；
  \item 两个函数同时变化时，用 Cauchy 的共享参数运输比值关系；
  \item 足够多个零点或重零点经过反复 Rolle，产生高阶导数见证。
\end{itemize}

因此区间定理不是“记住四个定理”，而是把\textbf{边界约束转译成内部见证}。

\begin{corebox}{中值点只保证存在，不给出可随意选择的固定坐标}
结论“存在 $\xi\in(a,b)$”是一个存在性输出。除非题设进一步唯一化或给出估计，不能把 $\xi$ 当成某个端点、对称点或自己挑选的数。所有回译都必须保留 $\xi$ 的合法区间和分母条件。
\end{corebox}

% ============================================================
\section{母模型：目标、辅助函数、资格、见证与回译}
% ============================================================

本册把中值定理和函数形状压缩为六个依次发生的动作：
\begin{processchain}
  \processstage{Target}
  \processstage{Auxiliary Function}
  \processstage{Qualification}
  \processstage{Boundary/Zero Design}
  \processstage{Witness}
  \processstage{Translate Back}
\end{processchain}

这里的箭头表示\textbf{证明和分析的控制顺序}：
\begin{enumerate}
  \item \kw{Target}：把待证结论写成 $\Phi(\xi)=0$ 或某个导数/符号目标；
  \item \kw{Auxiliary Function}：反推一个 $F$，使 $F^{(k)}$ 等于 $\Phi$ 的非零倍数；
  \item \kw{Qualification}：核对闭区间连续、开区间可导、分母非零、对数真数和积分因子；
  \item \kw{Boundary/Zero Design}：让 $F$ 端点等值，或制造足够多个不同零点/重零点；
  \item \kw{Witness}：应用介值、Rolle、Lagrange、Cauchy 或反复 Rolle，产生合法内部点；
  \item \kw{Translate Back}：展开 $F^{(k)}(\xi)=0$，用非零因子回译为原目标，并保留区间信息。
\end{enumerate}

\begin{center}
\begin{tikzpicture}[
  node distance=0.45cm and 0.48cm,
  box/.style={draw=iparaDeep,rounded corners=1.5mm,minimum width=3.0cm,minimum height=1.0cm,align=center,fill=iparaSoft,font=\small},
  arr/.style={-{Latex[length=2mm]},thick,draw=iparaDeep}
]
\node[box] (a) {目标表达式\\$\Phi(\xi)=0$};
\node[box,right=of a] (b) {反推辅助函数\\$F^{(k)}=m\Phi$};
\node[box,right=of b] (c) {资格审查\\连续/可导/非零};
\node[box,below=of c] (d) {制造约束\\端点等值/多零点};
\node[box,left=of d] (e) {中值见证\\$\xi$ 存在};
\node[box,left=of e] (f) {回译结论\\保留区间与条件};
\draw[arr] (a)--(b);
\draw[arr] (b)--(c);
\draw[arr] (c)--(d);
\draw[arr] (d)--(e);
\draw[arr] (e)--(f);
% Fallback 是外围反馈，不与主流程争用底部中心入射空间。
% a.south 下方已有 f，故 south 属于 blocked side；改从 a.west 的 free side 接入。
\coordinate (fallbackBottom) at ($(d.south east)+(0,-0.42)$);
\coordinate (fallbackLeft) at ($(a.west)+(-0.42,0)$);
\draw[arr,dashed,draw=iparaWarn]
  (d.south east) -- (fallbackBottom)
  -- node[midway,below,font=\scriptsize]{约束不足：停止套定理} (fallbackBottom -| fallbackLeft)
  -- (fallbackLeft) -- (a.west);
\end{tikzpicture}
\end{center}

函数形状题运行同一模型的逆向版本：先由 $f'$ 或 $f''$ 的符号得到局部/区间见证，再通过 Lagrange 把符号运输成函数值差异，最后确认极值或凹凸边界。

% ============================================================
\section{定理阶梯：连续、水平、割线与双函数运输}
% ============================================================

\subsection{闭区间连续：紧致性控制范围，连通性填补空隙}

点态连续只承诺：在\emph{每一个固定点}附近，输出可以被局部控制。当把定义域置于实数区间时，连续性会沿着两条不同的生成机制升级为全局结论：
\begin{enumerate}
  \item \kw{紧致性机制（闭且有界）}：输入既不能向无穷逃逸，也不会遗漏端点，使各点各自的局部控制能被升级为\textbf{全局统一控制}；
  \item \kw{连通性机制（区间无断裂）}：定义域是一整块没有缺口的连通集，连续映射把连通集映为连通集，从而\textbf{阻止输出产生跳跃或中空}。
\end{enumerate}

若 $f\in C[a,b]$，紧致性直接给出：
\begin{enumerate}
  \item \kw{有界性}：存在 $M>0$，使所有 $x\in[a,b]$ 都满足 $\lvert f(x)\rvert\le M$；
  \item \kw{最值存在}：存在 $x_m,x_M\in[a,b]$，使
  \[
  f(x_m)=\min_{[a,b]}f,
  \qquad
  f(x_M)=\max_{[a,b]}f.
  \]
\end{enumerate}

这两条不是“连续函数天然不会很大”，而是\kw{连续性 + 闭区间紧致性}共同作用的结果。证明骨架可以用 Topic02 的数列观点复原：若 $f$ 无界，就能选出 $x_n\in[a,b]$ 使 $\lvert f(x_n)\rvert>n$；闭区间中的数列存在收敛子列 $x_{n_k}\to x_*$，而连续性又迫使 $f(x_{n_k})\to f(x_*)$ 为有限值，矛盾。最值定理类似：先取函数值逼近上确界/下确界的点列，再用闭区间中的收敛子列和连续性把“逼近”升级成“真正取到”。

\begin{corebox}{闭区间连续的两大生成母模型}
\[
\begin{aligned}
\text{\bf 紧致性支链：}\quad &\text{局部控制}+\text{闭且有界（防逃逸/漏边界）}\Longrightarrow\text{有界、最值可达、一致连续} \\
\text{\bf 连通性支链：}\quad &\text{连续映射}+\text{区间连通（防跳跃断裂）}\Longrightarrow\text{介值定理、零点存在}
\end{aligned}
\]
介值性本质依赖区间的连通结构（在开区间、无穷区间上同样成立，不依赖紧致性）；而在闭区间上，紧致性给出确界点 $m,M$ 可达，连通性保证中间值被全覆盖，二者合力将像集锁定为紧闭区间 $[m,M]$。
\end{corebox}

\subsection{开区间有界：真正要封住的是两端逃逸}
闭区间只是最方便的“把输入关住”方式，不是有界性的唯一来源。若 $f$ 在 $(a,b)$ 上连续，并且两个单侧极限
\[
\lim_{x\to a^+}f(x),\qquad \lim_{x\to b^-}f(x)
\]
都存在且为\textbf{有限实数}，那么在两端各取一个足够小的邻域即可获得局部界；中间剩下的闭区间再由连续性给出界，因此 $f$ 在整个 $(a,b)$ 上有界。

若端点换成 $+\infty$ 或 $-\infty$，逻辑完全相同：要求相应无穷远极限为有限实数，再把中间有限部分截成闭区间。这里“有限”是不可删除的资格；例如 $1/x$ 在 $(0,1)$ 连续，但 $x\to0^+$ 时趋于 $+\infty$，所以无界。

\begin{warnbox}{有界不等于取到最值，连续也不等于有界}
$f(x)=x$ 在 $(0,1)$ 上有界，却没有最大值和最小值；$f(x)=1/x$ 在 $(0,1)$ 上连续，却无界。真正决定这些结论的不是一个孤立标签，而是“局部行为 + 定义域边界是否允许逃逸”。
\end{warnbox}

\subsection{介值定理与零点定理：连续性填补输出空隙}

若 $f$ 在 $[a,b]$ 上连续，则任意介于 $f(a)$ 与 $f(b)$ 的值 $\mu$ 都由某个 $c\in[a,b]$ 取得。若 $f(a)f(b)<0$，则存在 $c\in(a,b)$ 使 $f(c)=0$。

这不是“图像看起来会穿过横轴”的直觉，而是连续性阻止跳跃的存在性结论。若端点值恰好为零，见证点可能在端点；要求内部零点时必须明确端点非零或改变区间。

结合上一节的最值定理，若
\[
m=\min_{[a,b]}f,\qquad M=\max_{[a,b]}f,
\]
那么每个 $\mu\in[m,M]$ 都会被函数取到：取到最小值和最大值的两点之间应用介值定理即可。因此闭区间连续函数的像本身也是一个没有缺口的闭区间 $[m,M]$。

\subsection{一致连续：把每个点自己的输入尺度升级成全区间统一尺度}

普通连续在每个点 $x_0$ 允许使用不同的邻域尺度：
\[
\forall x_0\in[a,b],\ \forall\varepsilon>0,\ \exists\delta(x_0,\varepsilon)>0.
\]
一致连续则要求对给定 $\varepsilon$，存在一个与位置无关的统一尺度：
\[
\forall\varepsilon>0,\ \exists\delta>0,\ \forall x,y\in[a,b],
\quad \lvert x-y\rvert<\delta\Longrightarrow \lvert f(x)-f(y)\rvert<\varepsilon.
\]

若 $f\in C[a,b]$，则 $f$ 必一致连续（Heine--Cantor）。最能暴露机制的反证骨架是：若不能统一，就存在某个 $\varepsilon_0>0$ 和两列 $x_n,y_n\in[a,b]$，满足
\[
\lvert x_n-y_n\rvert\to0,
\qquad
\lvert f(x_n)-f(y_n)\rvert\ge\varepsilon_0.
\]
由闭区间可从 $x_n$ 取收敛子列 $x_{n_k}\to x_*$；因为两列距离趋零，同时有 $y_{n_k}\to x_*$。连续性便推出两列函数值都趋于 $f(x_*)$，与它们始终至少相差 $\varepsilon_0$ 矛盾。

\begin{warnbox}{点态连续不等于一致连续}
在非紧区间上，局部所需的 $\delta$ 可能随着位置不断缩小。例如 $f(x)=x^2$ 在 $\mathbb R$ 上处处连续，却不一致连续：同样小的输入差放到足够大的 $x$ 附近，会被放大成任意大的输出差。因此“一致连续”控制的是\kw{整片区域能否共用一把输入尺子}。
\end{warnbox}

\subsection{Rolle：端点等值制造内部水平切线}

若 $f$ 在 $[a,b]$ 连续、在 $(a,b)$ 可导且 $f(a)=f(b)$，则存在 $\xi\in(a,b)$ 使
\[
f'(\xi)=0.
\]

Rolle 是整个定理链的发动机：只要能把一个目标改写成“某个辅助函数两端相等”，就能获得内部导数为零的见证。端点等值是必须验证的边界约束，不是可以省略的形式。

\subsection{Lagrange：减去割线后恢复平均变化率}

若 $f$ 在 $[a,b]$ 连续、在 $(a,b)$ 可导，则取割线斜率
\[
K=\frac{f(b)-f(a)}{b-a},
\]
构造
\[
F(x)=f(x)-f(a)-K(x-a).
\]
此时 $F(a)=F(b)=0$。由 Rolle 得到
\[
F'(\xi)=f'(\xi)-K=0,
\]
即
\[
\boxed{f'(\xi)=\frac{f(b)-f(a)}{b-a}}.
\]

Lagrange 的机制不是“端点斜率等于某点斜率”，而是先减掉已知的线性趋势，把问题变成水平端点问题。

\subsubsection{有界导数在有限区间上给出利普希茨控制}
若 $I$ 是有限区间，且
\[
\lvert f'(x)\rvert\le K
\]
在区间内部成立，则对任意 $x,x_0\in I$，在连接两点的闭子区间上应用 Lagrange 得
\[
\lvert f(x)-f(x_0)\rvert\le K\lvert x-x_0\rvert.
\]
因此输出变化被输入距离线性控制；由于 $I$ 的长度有限，固定一个 $x_0$ 就得到 $f$ 的全区间有界性。这条结论真正使用的是
\[
\boxed{\text{导数有界}\Rightarrow\text{利普希茨控制},\qquad \text{有限输入跨度}\Rightarrow\text{输出有界}.}
\]

“有限区间”不能删除：$f(x)=x$ 在 $\mathbb R$ 上满足 $\lvert f'\rvert=1$，但函数本身无界。

\subsection{Cauchy：共同参数下的比值运输}

若 $f,g$ 在 $[a,b]$ 连续、在 $(a,b)$ 可导，则存在 $\xi$ 使
\[
[f(b)-f(a)]g'(\xi)=[g(b)-g(a)]f'(\xi).
\]
只有在进一步确保 $g'(\xi)\ne0$ 或对应差值非零时，才可把它改写为导数比。安全证明应优先保留交叉相乘形式，避免未经资格审查的除法。

可构造
\[
F(x)=[f(b)-f(a)][g(x)-g(a)]-[g(b)-g(a)][f(x)-f(a)],
\]
其两端均为零，再应用 Rolle。Cauchy 是 Lagrange 的双函数接口，不是“看到分式就自动洛必达”。

\subsection{L'Hospital：怎样把原商的极限运输到导数比？}

L'Hospital 法则真正解决的不是“分式不好算”，而是下面这个更窄的问题：当分子、分母同时趋于 $0$，或同时以无界尺度增长时，\textbf{原商的相对尺度能否由局部变化率之比稳定决定？}

以有限点 $x\to a$ 的 $0/0$ 型为母型。设 $f,g$ 在 $a$ 的某个去心邻域内可导，$g'(x)\ne0$，并且
\[
 f(x)\to0,\qquad g(x)\to0,
\]
若导数比极限
\[
\lim_{x\to a}\frac{f'(x)}{g'(x)}=L
\]
存在（允许为扩展极限），则在相应资格下有
\[
\boxed{\lim_{x\to a}\frac{f(x)}{g(x)}=L.}
\]
单侧极限与无穷远版本遵循同一机制，但必须把区间和趋近方向改成对应的合法尾部。

\begin{corebox}{为什么 Cauchy 中值定理能生成 L'Hospital}
先看 $0/0$ 型。对靠近 $a$ 的 $x$，若把端点值按极限补成 $f(a)=g(a)=0$，并满足相应连续资格，则在 $a$ 与 $x$ 之间由 Cauchy 中值定理存在 $\xi_x$，使
\[
\frac{f(x)-f(a)}{g(x)-g(a)}
=
\frac{f'(\xi_x)}{g'(\xi_x)}.
\]
左边就是 $f(x)/g(x)$。当 $x\to a$ 时，中值点 $\xi_x\to a$；若导数比对所有合法趋近都稳定到 $L$，原商也只能稳定到 $L$。

这解释了法则的方向：\textbf{不是“分子分母分别求导以后数值不变”，而是 Cauchy 中值定理在一个不断收缩的区间内，把端点差值比运输成某个内部点的导数比。}
\end{corebox}

$\infty/\infty$ 型不能简单写成“在极限点补零”，但仍使用同一比值运输思想：选取一个固定远端基点，用 Cauchy 中值定理比较两个端点的增量，再利用 $f,g$ 的无界主导项压过固定基点贡献。完整技术形式因有限点、单侧和无穷远版本略有差异；考研使用时最重要的是下面的统一资格链。

\begin{methodbox}{L'Hospital 的资格链}
每次调用都依次核对：
\begin{enumerate}
  \item \kw{型别}：当前商确实是 $0/0$ 或 $\infty/\infty$；其他五类未定式必须先在 Topic02 标准化；
  \item \kw{区间}：分子分母在对应去心邻域或远端尾部内可导，而不是只在目标点“形式上能求导”；
  \item \kw{分母变化率}：$g'(x)\ne0$ 在所需邻域内成立；
  \item \kw{导数比}：$f'(x)/g'(x)$ 的极限存在，或至少已经进入可继续分析的合法结构；
  \item \kw{收益}：一次求导后表示确实更简单，或至少暴露了原式所缺的尺度信息。
\end{enumerate}
\end{methodbox}

连续多次使用时，\textbf{不能把“允许使用一次”理解成“以后可以一直求导”}。第一次求导后得到的是一个新商；若想再次使用，必须重新确认这个新商仍为 $0/0$ 或 $\infty/\infty$，并重新检查可导性、分母导数与复杂度。

\begin{warnbox}{L'Hospital 的典型失败不是“函数太复杂”，而是资格或表示已经坏掉}
\begin{itemize}
  \item 绝对值、分段、取整等在规则切换点可能破坏去心邻域可导性；必要时先拆左右极限。
  \item 强振荡项求导后可能把振荡放大，例如出现 $\sin(1/x)$、$\cos(1/x)$ 与负幂组合；此时夹逼或归结原则往往更接近问题本体。
  \item 根式、分数幂和对数必须先维护真实定义域；不能为了“凑可导”越过非法分支。
  \item 若求导后已经不是未定式，应直接求新商极限，不能因为结果尚未一眼可见就再次机械求导。
  \item 若一阶主部相消、参数需要配阶，Topic03 的有限 Taylor 通常比连续 L'Hospital 更能暴露最低非零项。
\end{itemize}
\end{warnbox}

因此 L'Hospital 与 Taylor、等价无穷小的区别不是“哪个更高级”：L'Hospital 用\textbf{变化率比}搬运商极限；Taylor 显式保留\textbf{若干阶局部系数与余项}；等价无穷小只保留\textbf{稳定主部}。目标只需要哪一级信息，就用哪一级机制。

\subsection{高阶 Rolle：零点数量转为高阶导数见证}

若 $F$ 有 $n+1$ 个互不相同的零点，并具备相应阶数的连续/可导条件，则反复应用 Rolle 可得某个 $\xi$ 使
\[
F^{(n)}(\xi)=0.
\]
重零点按其导数消失阶数计数。例如 $F(a)=F'(a)=0$ 等价于在反复 Rolle 中提供两个零点条件。中间每一次得到的点都位于相邻零点之间，不能把不同层的中值点混为同一个点。

% ============================================================
\section{证明题逆向工程：从目标表达式反推辅助函数}
% ============================================================

\subsection{先把结论标准化}

把“存在 $\xi$ 使某式成立”写成
\[
\Phi(\xi)=0.
\]
然后问：能否找到 $F$ 与在区间内不为零的 $m$，使
\[
F'(x)=m(x)\Phi(x)
\]
或更高阶地 $F^{(k)}(x)=m(x)\Phi(x)$？这一步把猜辅助函数变成了可检查的逆向设计。

\subsection{结构反推表}

\begin{handbooklongtable}{L{0.20\linewidth}L{0.28\linewidth}Y}
\toprule
\textbf{目标导数结构} & \textbf{辅助函数} & \textbf{必须核对的资格}\\
\midrule
$f'(x)=K$ & $f(x)-Kx$ 或减去割线 & 端点等值、区间连续/可导\\
$A'B+AB'$ & $AB$ & 因子在区间内有定义\\
$A'B-AB'$ & $A/B$ & $B$ 在区间内不为零\\
$A'/A\pm B'/B$ & $\ln\lvert A\rvert\pm\ln\lvert B\rvert$ & $A,B$ 不为零且符号区间固定\\
$y'+py=q$ & 积分因子乘积减变限积分 & $p,q$ 连续，指数与积分合法\\
$F^{(n)}(x)=K$ & 减去满足插值条件的低次多项式 & 足够多个零点/重零点\\
\bottomrule
\end{handbooklongtable}

表格提供的是机制接口，不是无条件题型口诀。每次使用仍须先完成 Qualification 和 Boundary\allowbreak/Zero Design 两步。

\subsection{乘积、商与对数结构}

若目标含 $A'B+AB'$，逆用乘积法则取 $F=AB$；若含 $A'B-AB'$ 且 $B\ne0$，取 $F=A/B$。若含 $A'/A$，对数辅助函数可把乘法结构压成加法，但真数非零与符号保持是不可删除的资格。

例如要证明存在 $\xi$ 使
\[
f'(\xi)[g(\xi)-g(b)]+g'(\xi)[f(\xi)-f(a)]=0,
\]
取
\[
F(x)=[f(x)-f(a)][g(x)-g(b)].
\]
有 $F(a)=F(b)=0$，Rolle 给出 $F'(\xi)=0$，展开即得目标。真正的设计动作是让两个因子分别在相反端点消失。

\subsection{积分因子与变上限函数只作局部工具}

对一阶线性结构 $y'+p(x)y=q(x)$，令
\[
\mu(x)=\exp\left(\int p(x)\,dx\right)
\]
使 $[\mu y]'=\mu(y'+py)$。非齐次目标可以再减去一个变上限积分，使辅助函数的导数等于方程残差的非零倍数。

但“总能定义 $F(x)=\int\Phi$”不自动产生中值点；还必须另外证明 $F$ 在两个点等值或拥有足够零点。积分只解决导数匹配，不替代边界设计。

% ============================================================
\section{函数形状：把导数符号运输成区间性质}
% ============================================================

\subsection{单调性是 Lagrange 的全区间回译}

若 $f$ 在区间 $I$ 上连续、在内部可导，且对任意 $x$ 有 $f'(x)>0$，任取 $x_1<x_2$，由 Lagrange 存在 $\xi\in(x_1,x_2)$ 使
\[
f(x_2)-f(x_1)=f'(\xi)(x_2-x_1)>0.
\]
所以 $f$ 严格递增；$f'<0$ 同理严格递减；$f'=0$ 则函数为常数。若只知道 $f'\ge0$，同一中值定理直接给出 $f$ 单调不减。

反方向也要分层：若 $f$ 在区间内可导且单调不减，则任一点的差商都非负，极限只能满足
\[
\boxed{f'(x)\ge0}.
\]
严格递增仍只能推出 $f'\ge0$，不能推出 $f'>0$。最小反例是
\[
f(x)=x^3,
\]
它在 $\mathbb R$ 上严格递增，但 $f'(0)=0$。因此应记成
\[
\boxed{f'>0\Rightarrow\text{严格增},\qquad \text{严格增}+\text{可导}\Rightarrow f'\ge0,}
\]
而不是双向写成 $f'>0$。

若 $f'\ge0$ 但在某一整段上恒为 $0$，函数会出现平台，只能保证单调不减；想从 $f'\ge0$ 升级到严格递增，还要排除任何非退化子区间上 $f'\equiv0$ 的情况。

\subsection{极值：先列候选，再检查符号变化}

内部极值点若可导，必满足 $f'(a)=0$；不可导尖点也可能是极值。因此候选集至少包括：
\begin{itemize}
  \item 驻点 $f'(a)=0$；
  \item $f'$ 不存在但函数有定义的点；
  \item 定义域端点（若问闭区间绝对极值）。
\end{itemize}

候选不等于结论。最稳的确认方法是检查 $f'$ 在两侧的符号：负到正为极小，正到负为极大，符号不变则不是极值。若 $f'(a)=0$ 且 $f''(a)>0$ 或 $<0$，二阶充分条件可快速确认；$f''(a)=0$ 时必须回到符号变化或更高阶第一非零导数。

\subsection{凹凸与拐点：二阶符号也要看区间}

在二阶可导区间内，$f''>0$ 表示斜率递增、图像凹向上；$f''<0$ 表示斜率递减、图像凹向下。拐点的核心是凹凸性在点两侧发生改变，不是单独看到 $f''(a)=0$。

若 $f''(a)=0$ 且 $f^{(3)}(a)\ne0$，在相应连续性资格下通常可确认拐点；若二阶导数在点不存在，仍可通过两侧凹凸定义判断。不可导点可以是拐点，但必须检查函数连续和凹凸真正改变。

\subsection{三点割线序与导函数单调：严格性要分两层恢复}

二阶导数只是判断凸性的一条充分路径，不是唯一入口。若 $f$ 在 $(a,b)$ 可导，则下面两件事等价：
\begin{enumerate}
  \item $f'$ 在 $(a,b)$ 严格递增；
  \item 任意 $x_1<x_2<x_3$ 都满足相邻割线斜率严格递增：
  \[
  \frac{f(x_2)-f(x_1)}{x_2-x_1}
  <
  \frac{f(x_3)-f(x_2)}{x_3-x_2}.
  \]
\end{enumerate}

由 1 推 2 很直接：分别在 $[x_1,x_2]$ 与 $[x_2,x_3]$ 使用 Lagrange，得到
$\xi_1\in(x_1,x_2)$、$\xi_2\in(x_2,x_3)$。因为 $\xi_1<\xi_2$，两条割线斜率分别等于
$f'(\xi_1)$ 与 $f'(\xi_2)$，严格次序随即成立。

反向不能把严格割线不等式直接``取极限''成严格导数不等式。任取 $u<v$，再取
$x<u<v<y$，由三点条件两次得到
\[
\frac{f(u)-f(x)}{u-x}
<
\frac{f(v)-f(u)}{v-u}
<
\frac{f(y)-f(v)}{y-v}.
\]
令 $x\to u^-$、$y\to v^+$，只能推出
\[
f'(u)\le
\frac{f(v)-f(u)}{v-u}
\le f'(v),
\]
所以第一层结论只是 $f'$ 非减。

严格性来自第二层的\textbf{平台排除}。若某个 $u<v$ 满足 $f'(u)=f'(v)$，非减性迫使
$f'$ 在 $[u,v]$ 上恒定；由 Lagrange，$f$ 在该段为仿射函数。于是在 $(u,v)$ 内任取三个点，相邻割线斜率相等，与严格三点条件矛盾。因此任意 $u<v$ 都有 $f'(u)<f'(v)$。

\begin{warnbox}{严格不等式经过极限可能降级}
若 $A_x<B<C_y$，即使 $A_x\to A$、$C_y\to C$，通常也只能得到 $A\le B\le C$。本题必须先得到导函数非减，再另用``若端点导数相等则整段为平台''恢复严格性。直接写``令端点趋近，所以 $f'(u)<f'(v)$''缺少关键证明。
\end{warnbox}

\subsection{零点重数与符号穿越}

若
\[
f(x)={(x-a)}^n g(x),
\qquad g(a)\ne0,
\]
则 $n$ 的奇偶决定函数在 $a$ 两侧是否变号：偶数阶通常保号并产生极值型接触，奇数阶变号；当奇数阶至少为三时，常伴随凹凸翻转。这里“通常/常伴随”不能替代直接检查 $g$ 的符号与凹凸。

对于多项式，$f'$ 的奇数重根会导致 $f'$ 变号并给出极值横坐标；$f''$ 的奇数重根会导致凹凸翻转并给出拐点横坐标。这个技巧依赖多项式/连续因子结构，不能无条件推广到带振荡或不规则因子的函数。

\subsection{参数曲线的渐近线：先穷举参数逃逸事件，再判断平面中的逃逸方向}

对参数曲线
\[
x=x(t),\qquad y=y(t),\qquad t\in D,
\]
渐近线分析的第一步不是直接计算 $y/x$，而是先找所有可能让曲线“走向无穷远”的\kw{参数临界事件}：
\begin{itemize}
  \item 参数定义域的有限端点与内部断点；
  \item 每个断点允许的左右单侧趋近；
  \item 若定义域无界，还包括 $t\to+\infty$ 与 $t\to-\infty$。
\end{itemize}
对每一个事件分别追踪
\[
L_x=\lim x(t),\qquad L_y=\lim y(t).
\]
然后按平面中的实际逃逸方向分类：
\begin{enumerate}
  \item $L_x=\pm\infty$、$L_y=B\in\mathbb R$：水平渐近线 $y=B$；
  \item $L_x=A\in\mathbb R$、$L_y=\pm\infty$：垂直渐近线 $x=A$；
  \item $\lvert x(t)\rvert\to\infty$ 且 $\lvert y(t)\rvert\to\infty$：尝试求
  \[
  k=\lim\frac{y(t)}{x(t)},\qquad
  b=\lim\bigl(y(t)-kx(t)\bigr).
  \]
  若 $k,b$ 都为有限实数，则 $y=kx+b$ 是该分支的仿射渐近线；$k=0$ 时退化为水平情形。
  \item $L_x,L_y$ 都有限：该参数事件只对应有限端点/缺口，不产生“走向无穷远”的渐近线。
\end{enumerate}

\begin{corebox}{参数曲线渐近线的生成模型}
\[
\boxed{\text{参数临界事件全集}\to(L_x,L_y)\to\text{逃逸方向}\to\text{水平/垂直/仿射渐近线}}
\]
关键不变量是\textbf{每个合法单侧事件都必须核销}。同一个有限参数断点的左右两侧可能生成不同渐近行为；只检查 $t\to\pm\infty$ 会系统性漏解。这里本册负责“参数事件枚举与平面几何组装”，单侧/无穷极限计算直接调用 \kw{Topic02} 极限发动机。
\end{corebox}

% ============================================================
\section{贯穿母例：三次函数的一条完整区间轨迹}
% ============================================================

令
\[
f(x)=x^3-3x.
\]
这是一个同时暴露候选点、符号运输、形状边界和辅助函数构造的母例。

\subsection{从导数符号到形状}

\[
f'(x)=3(x^2-1)=3(x-1)(x+1).
\]
因此 $f'>0$ 于 $(-\infty,-1)$ 与 $(1,\infty)$，$f'<0$ 于 $(-1,1)$。Lagrange 回译出：函数先增、后减、再增；$x=-1$ 是极大点，$x=1$ 是极小点。

再看
\[
f''(x)=6x.
\]
二阶导数在 $0$ 左负右正，故凹凸发生改变，$0$ 是拐点。注意 $f'(0)=-3\ne0$，拐点不要求驻点。

\subsection{同一个对象上的中值见证}

在任意 $a<b$ 上，令
\[
K=\frac{f(b)-f(a)}{b-a},
\qquad F(x)=f(x)-f(a)-K(x-a).
\]
则 $F(a)=F(b)=0$，所以存在 $\xi\in(a,b)$ 使 $f'(\xi)=K$。这说明“割线斜率被某条切线实现”并不依赖三次函数的特殊可解性，而是来自辅助函数与 Rolle 的通用结构。

\subsection{典型错误路径}

若只解 $f'(x)=0$ 并写“得到两个极值”，就漏掉了符号变化验证；若只解 $f''(x)=0$ 并写“得到拐点”，就漏掉了凹凸变化验证；若在中值证明中直接指定 $\xi=(a+b)/2$，就把存在性见证错误地固定成坐标。三处错误分别丢失了 Candidate/Confirm、区间符号和中值点存在性边界。

% ============================================================
\section{概念边界：真正判据、题目信号与错误后果}
% ============================================================

\begin{handbooklongtable}{L{0.13\linewidth}C{0.035\linewidth}L{0.13\linewidth}YY}
\toprule
\textbf{概念 A} & & \textbf{概念 B} & \textbf{为什么容易混 / 真正判据与题目信号} & \textbf{混淆后果}\\
\midrule
介值定理 & $\neq$ & Rolle & 介值只需连续并产生函数值；Rolle 还需内部可导和端点等值，产生导数零点。 & 用连续性直接推出内部水平切线\\
Rolle & $\neq$ & Lagrange & Lagrange 先减割线再用 Rolle；目标是平均变化率，不是只能证明零导数。 & 忘记构造割线修正\\
Lagrange & $\neq$ & Cauchy & Cauchy 同时运输两个函数的增量；比值形式还需分母资格。 & 未审查分母就相除\\
端点等值 & $\neq$ & 端点同号 & Rolle 需要 $F(a)=F(b)$；零点定理需要跨值/异号，二者用途不同。 & 把零点存在误当水平切线存在\\
零点 & $\neq$ & 重零点 & $F(a)=0$ 只计一个约束；还需 $F'(a)=0$ 等条件才可在反复 Rolle 中多次计数。 & 高阶导数见证少算或多算\\
驻点 & $\neq$ & 极值点 & $f'=0$ 只产生候选；一阶导数变号或充分条件才确认。 & 把平台拐点误判为极值\\
二阶导数为零 & $\neq$ & 拐点 & 拐点要求凹凸改变；$f''=0$ 只是候选，且不可导点也可能是拐点。 & 看到二阶零点就宣布拐点\\
严格割线序 & $\neq$ & 取极限后仍严格 & 严格不等式取极限一般只保留非严格次序；需先证 $f'$ 非减，再排除常值平台。 & 省略严格单调性的关键一步\\
局部导数符号 & $\neq$ & 全区间单调 & 单调结论需要对任意子区间调用 Lagrange 或其他全区间论证；单点符号没有全局信息。 & 从一个点的斜率外推整段函数\\
中值点存在 & $\neq$ & 中值点坐标已知 & 定理只给某个 $\xi$ 存在及其区间，通常不允许把它固定成对称点。 & 在证明中偷换见证点\\
直接积分 & $\neq$ & 已完成 Rolle 证明 & $F'=\Phi$ 只解决导数匹配；仍需端点等值或多零点。 & 以“找到原函数”替代存在性条件\\
\bottomrule
\end{handbooklongtable}

% ============================================================
\section{Topic04 的问题求解协议}
% ============================================================

面对中值定理证明、零点、单调极值或函数形状问题，按下列顺序维护状态：
\[
\boxed{
\begin{gathered}
\text{目标标准化}\to\text{资格审查}\to\text{辅助函数}\to\text{边界/零点设计}\\
\to\text{定理见证}\to\text{回译}\to\text{区间/符号验证}
\end{gathered}}
\]

\subsection{Step 1：Target}

把待证结论改写为 $\Phi(\xi)=0$、$f'(\xi)=K$、$F^{(n)}(\xi)=0$ 或函数值符号差。若问形状，先标出候选点与目标区间。

\subsection{Step 2：Qualification}

逐项检查闭区间连续、开区间可导、分母非零、对数真数、积分因子和端点是否属于定义域。资格失败时立即换定理或缩小合法区间。

\subsection{Step 3：Auxiliary Function}

逆用和积商链、割线修正、对数导数、积分因子或插值多项式，使某阶导数成为目标的非零倍数。辅助函数越短且端点越容易验证越好。

\subsection{Step 4：Boundary / Zero Design}

检查 $F(a)=F(b)$，或列出不同零点与重零点的区间顺序。需要多个中值点时分割区间；不能默认多次定理给出同一个点。

\subsection{Step 5：Witness and Translate Back}

明确引用介值、Rolle、Lagrange、Cauchy 或反复 Rolle，写出 $\xi$ 的所在区间，再展开 $F^{(k)}(\xi)=0$ 并除以已核对的非零因子。

\subsection{Step 6：Shape Verification}

做函数形状题时，按“候选集 → 一阶符号 → 二阶符号 → 端点比较”的顺序。独立检查包括：极值左右符号是否翻转、拐点凹凸是否改变、闭区间绝对极值是否比较端点、根的数量是否与 Rolle 下界相容。

% ============================================================
\section{高频失败路径：第一次偏离发生在哪里？}
% ============================================================

\begin{enumerate}
  \item \kw{资格失败}：未检查闭区间连续/开区间可导，就直接写“由中值定理”。
  \item \kw{目标失败}：没有先把结论标准化为导数或零点结构，盲目猜辅助函数。
  \item \kw{边界失败}：辅助函数导数匹配了目标，却没有制造端点等值或足够零点。
  \item \kw{分母失败}：Cauchy、商函数、对数导数或积分因子中分母可能为零，却直接相除。
  \item \kw{洛必达资格失败}：未先确认 $0/0$ 或 $\infty/\infty$，或连续求导时没有把每一次新商重新验型、验可导性与验收益。
  \item \kw{见证失败}：把存在的中值点固定成中点，或把多次应用得到的点误认为同一点。
  \item \kw{候选失败}：由 $f'=0$、$f''=0$ 或高阶导数为零直接宣布极值/拐点。
  \item \kw{严格性失败}：把严格不等式取极限后仍写成严格，未用平台反证恢复严格性。
  \item \kw{区间失败}：用一个点的导数符号替代整个区间的符号，或忘记闭区间端点比较。
  \item \kw{Owner 失败}：把插值余项、无限 Taylor、变限积分或完整 ODE 解法复制进本册。
\end{enumerate}

% ============================================================
\section{传统章节与 Source Corpus 映射}
% ============================================================

\begin{handbooklongtable}{L{0.16\linewidth}L{0.32\linewidth}Y}
\toprule
\textbf{Source / 传统内容} & \textbf{进入本册} & \textbf{分流与拒绝复制}\\
\midrule
I-07 定理链 & 介值/零点、Rolle、Lagrange、Cauchy、高阶 Rolle 的资格与生成关系 & 题目识别和首步动作进入同目录训练 Markdown\\
I-02 函数极限中的 L'Hospital & Cauchy 中值定理到商极限的比值运输机制、$0/0$ 与 $\infty/\infty$ 资格链、连续使用时的逐次重检 & 七类未定式标准化与尺度诊断归 Topic02；具体题面触发、换路和退出动作进入同目录训练 Markdown\\
I-07 证明逆向工程 & 目标标准化、辅助函数反推、端点等值/零点设计与回译 & 只保留方法机制，不复制成题型清单\\
I-07 K 值法 & 割线修正、减去低次多项式、重零点计数 & 插值余项与误差估计移交 H-B02\\
I-07 乘积/商/对数/积分因子 & 作为辅助函数构造的局部结构接口 & 乘积/商本体归 Topic03；变限积分正则性归 Topic05/H-B03；ODE 解族归 Topic12\\
I-07 零点重数与反复 Rolle & 高阶导数见证、根数量下界和重根边界 & 多中值点插点法归 H-B02\\
I-06 函数形状 & 单调、极值、凹凸、拐点、零点重数与图像读法 & 曲率/局部二阶模型归 Topic03；定义域/函数对象归 Topic01\\
I-04 参数曲线渐近线 & 参数临界事件全集、单侧极限与水平/垂直/仿射渐近线分类 & 具体题面扫描顺序进入同目录训练 Markdown\\
I-07.1 多中值点与中值点参数 & 只保留“分割区间保证不同中值点”的边界提示 & 主体归 H-B02，不在 Topic04 重写中值点参数渐近\\
I-07.2 插值余项与导数估计 & 只保留“高阶 Rolle 接收插值条件”的接口 & Lagrange/Hermite 余项、Landau\textendash{}Kolmogorov 和误差界归 H-B02\\
\bottomrule
\end{handbooklongtable}

% ============================================================
\section{一页压缩：忘掉公式后怎样重建？}
% ============================================================

\begin{corebox}{一句话}
把目标表达式反推成辅助函数的导数，先验证区间资格，再制造端点等值或足够零点；中值定理只产生存在性见证，最后必须把见证回译并检查区间符号。
\end{corebox}

\subsection{母图}
\begin{processchain}
  \processstage{目标}
  \processstage{辅助函数}
  \processstage{资格}
  \processstage{端点/零点}
  \processstage{中值见证}
  \processstage{回译与形状}
\end{processchain}

\subsection{三个生成原则}
\begin{enumerate}
  \item \kw{运输原则}：介值运输函数值，Rolle 运输端点等值到内部零导数，Lagrange 运输割线斜率，Cauchy 运输双函数比值。
  \item \kw{设计原则}：辅助函数由目标导数结构与边界条件共同反推；只匹配导数而不制造边界约束是不完整的证明。
  \item \kw{候选原则}：驻点、二阶零点和中值点都是候选/存在性输出，必须用符号、区间和资格完成确认。
\end{enumerate}

\subsection{重建问题}
\begin{enumerate}
  \item 目标能否写成 $\Phi(\xi)=0$ 或某个导数等式？
  \item 需要哪个定理的资格：连续、可导、端点等值、非零分母还是多零点？
  \item 哪个辅助函数的导数最接近目标？它的端点值怎样验证？
  \item 得到的中值点位于哪个开区间，是否可能与另一个点混淆？
  \item $f'$ 或 $f''$ 的符号在区间上怎样变化？候选点是否真的确认？
  \item 若严格不等式经过了极限，严格性来自哪里？是否还需排除常值平台？
  \item 本题是否已经越界到 H-B02 的插值余项、Topic05 的积分接口或 Topic12 的 ODE 解族？
\end{enumerate}

\begin{center}\rule{0.5\linewidth}{0.5pt}\end{center}

\appendix
\section{中值定理与函数形状 Coverage 锚点}

\subsection{定理资格速查}

\begin{handbooktable}{L{0.16\linewidth}L{0.36\linewidth}Y}
\toprule
定理 & 资格 & 输出\\
\midrule
介值定理 & $[a,b]$ 连续 & 介于端点值之间的函数值存在\\
零点定理 & $[a,b]$ 连续且端点异号 & 内部零点存在\\
Rolle & 闭区间连续、开区间可导、端点等值 & $f'(\xi)=0$\\
Lagrange & 闭区间连续、开区间可导 & $f'(\xi)$ 等于割线斜率\\
Cauchy & 两函数均闭区间连续、开区间可导 & 交叉相乘的导数关系\\
高阶 Rolle & 足够阶光滑、零点/重零点条件 & 高阶导数零点存在\\
\bottomrule
\end{handbooktable}

\subsection{形状判定 Coverage}

按考纲查漏，本册覆盖：闭区间连续函数性质、介值与零点、Rolle\allowbreak/Lagrange\allowbreak/Cauchy 中值定理、由 Cauchy 生成的 L'Hospital 法则及其资格链、单调性、极值、最大最小值、凹凸与拐点、函数图像和导数根的数量。Lagrange\allowbreak/Hermite 插值余项、导数估计和中值点参数渐近由 H-B02 负责；七类未定式标准化归 Topic02；单一问题族的题面识别与路径选择进入同目录训练 Markdown。

\end{document}
